
% Introdução
%--------------------------------------
\chapter{Introdução}
A segurança das infraestruturas de rede assume um papel cada vez mais relevante no contexto atual, caracterizado por um aumento significativo na sofisticação e frequência das ameaças informáticas. A monitorização contínua do tráfego de rede tornou-se um elemento essencial para a deteção precoce de comportamentos anómalos ou maliciosos e para a mitigação de incidentes de segurança, contribuindo para a preservação da confidencialidade, integridade e disponibilidade da informação.

Os \textbf{Sistemas de Deteção de Intrusões (IDS)} e os \textbf{Sistemas de Prevenção de Intrusões (IPS)} surgem como mecanismos fundamentais neste domínio. Um \textbf{IDS} opera de forma passiva, através da análise do tráfego de rede, com base na comparação com padrões previamente definidos ou na identificação de comportamentos anómalos. Sempre que deteta indícios de atividade maliciosa, emite alertas que permitem a adoção de medidas corretivas adequadas. Um \textbf{IPS}, por sua vez, introduz uma vertente ativa no processo de proteção e não se limita à deteção de ameaças, ou seja, procede ao bloqueio, à rejeição ou à mitigação de comunicações maliciosas em tempo real, de modo a salvaguardar a integridade, a confidencialidade e a disponibilidade dos sistemas.

Deste modo, estas soluções integram-se numa estratégia de defesa, na qual assumem funções de monitorização e controlo que complementam outras medidas de segurança, como \textit{firewalls} e políticas de controlo de acesso. Assim, no âmbito da Unidade Curricular de Segurança, revela-se fundamental a compreensão do seu funcionamento nas vertentes teórica e prática.
